\section{Simulation studies}\label{sec:sim}

\subsection{Gaussian stress design and response mechanisms}

The first design isolates the missingness comparison with oracle nuisance functions, then examines the effect of estimating those functions. There are no fully observed baseline covariates $W$. The full-data law is specified retrospectively:
\begin{align*}
A&\sim\mathrm{Bernoulli}(0.5),\\
Y\mid A=0&\sim\mathrm{Beta}(2,4),\qquad
Y\mid A=1\sim\mathrm{Beta}(4,2),\\
L_1\mid A,Y&\sim\mathrm{Bernoulli}\{\expit(-0.6+0.5A+0.25Y+0.1AY)\},\\
L_2\mid A,Y,L_1&\sim N(A+Y+2.5L_1Y,1.25^2).
\end{align*}
Refactorization of this law yields $e(V)$ and $\mu_a(V)$ by numerical quadrature. The ATE is $\psi=0.2558505161$ to ten decimal places; all summaries use the unrounded computed value. Stage-one response is
\[
\pi_1=\expit(2+0.1A+0.1Y).
\]
The two stage-two mechanisms are
\[
\pi_2=\begin{cases}
\expit(-1.1+0.1A+0.1Y),&\text{common validity},\\
\expit(-1.9+0.1A+0.1Y+2.2L_1),&\text{sequential MAR only}.
\end{cases}
\]
The first permits an efficiency comparison and the second illustrates identification failure after coarsening. The oracle experiment uses 2,000 samples of size 2,500 per mechanism. All three estimators are means of their influence-function pseudo-outcomes, with normal intervals based on the empirical variance within each sample. The coarsened oracle procedure uses the true causal functions and $Q_0$ but weights by $1/\rho(Z_0)$; the supplement also reports the alternative augmentation $Q_0^{\cc}$. Neither oracle version is a feasible causal estimator outside common validity.

This Gaussian design does not satisfy the uniform treatment-positivity condition in \cref{ass:causal}. The unbounded support of $L_2$ permits treatment probabilities arbitrarily close to zero or one. The weaker moment conditions $\E(1/e)=4.04$ and $\E\{1/(1-e)\}=5.65$ nevertheless give a finite full-data efficiency bound because $Y$ is bounded. Deterministic quadrature gives $\Var(G)=0.27924$ and $\E(G^4)=105.28$. These features make this a stress design for learning the causal nuisance functions; it is not an empirical verification of \cref{thm:asymptotics}.

\subsection{Oracle comparison}

\begin{table}[htbp]
\centering
\begin{threeparttable}
\caption{Oracle simulation in the Gaussian stress design ($n=2{,}500$, 2,000 replicates per scenario).}
\label{tab:main-results}
\small
\setlength{\tabcolsep}{4pt}
\begin{tabular}{@{}llrrrr@{}}
\toprule
Scenario & Estimator & Bias & SD & RMSE & Coverage \\
\midrule
Common validity & Full data & $-0.00005$ & $0.01043$ & $0.01043$ & 94.95\% \\
 & Coarsened & $-0.00004$ & $0.01775$ & $0.01775$ & 95.10\% \\
 & Sequential & $-0.00003$ & $0.01723$ & $0.01722$ & 95.55\% \\
\addlinespace
Sequential MAR only & Full data & $0.00013$ & $0.01052$ & $0.01051$ & 95.05\% \\
 & Coarsened & $-0.03301$ & $0.01848$ & $0.03783$ & 45.70\% \\
 & Sequential & $-0.00013$ & $0.01626$ & $0.01626$ & 94.65\% \\
\addlinespace
\bottomrule
\end{tabular}
\begin{tablenotes}[flushleft]\footnotesize
\item The target ATE is $0.2558505161$. The coarsened column uses true $G$, $Q_0$, and $\rho$. Bias and coverage Monte Carlo standard errors are recorded in the supplement; coverage MCSE is at most 1.12 percentage points.
\end{tablenotes}
\end{threeparttable}
\end{table}

The centered empirical variance reduction under common validity was 5.83\%, with paired-bootstrap Monte Carlo 95\% interval 4.03\%--7.57\% (10,000 resamples). This statistic uses deviations from the replicate means, rather than squared deviations from the true ATE, which would estimate a mean-squared-error reduction.

Deterministic product quadrature gives influence-function variances 0.79524 and 0.76224 for the coarsened and sequential procedures, respectively. Their difference is 0.03300, corresponding to a 4.15\% reduction. Doubling the quadrature resolutions changes each variance by less than $10^{-8}$, and the numerical residual in \cref{eq:variance-gap} is below $10^{-12}$. These are converged numerical evaluations, not symbolic integrations. The 4.15\% value lies within the simulation interval; the interval's width also shows why a finite Monte Carlo variance ratio should not be equated with the efficiency bound.

In the sequential-only scenario, the coarsened oracle bias was $-0.03301$ with 45.70\% coverage, whereas the sequential bias was $-0.00013$ with 94.65\% coverage. Quadrature gives population coarsened bias $-0.03282$. Replacing $Q_0$ by $Q_0^{\cc}$ gives simulated bias $-0.03302$; the two population biases coincide, as shown above. This is an identification comparison.

\begin{figure}[htbp]
\centering
\includegraphics[width=\textwidth]{estimator_comparison.pdf}
\caption{Oracle estimates from the same replicates used in \cref{tab:main-results}. Boxes show quartiles, whiskers extend to 1.5 times the interquartile range, and points beyond the whiskers are shown. Dashed lines mark the true ATE.}
\label{fig:main-boxplot}
\end{figure}

The population pattern proportions obtained by quadrature are 23.98\% complete, 65.09\% $L_1$ only, and 10.93\% neither under common validity. The corresponding sequential-only proportions are 31.18\%, 57.89\%, and 10.93\%.

\subsection{Sensitivity to the second-stage response mechanism}

The sensitivity experiment keeps the Gaussian full-data law and uses $\pi_2=\expit(b+0.1A+0.1Y+\gamma L_1)$. Under common validity, $\gamma=0$ and the intercept $b$ is calibrated to population complete-case proportions 0.5117, 0.2400, and 0.1161. In the identification experiment, $\gamma$ is 0.75, 1.50, or 2.25 and $b$ is calibrated separately to maintain a population complete-case proportion of 0.24. Each configuration uses 5,000 samples of size 2,500. The intercepts, empirical response shares, biases, coverage, and uncertainty in the variance reduction are reported below. Calibration uses deterministic integration; reported empirical shares need not equal the population targets exactly.

\begin{table}[htbp]
\centering
\begin{threeparttable}
\caption{Efficiency sensitivity under common validity (5,000 replicates per row).}
\label{tab:sensitivity-efficiency}
\small
\setlength{\tabcolsep}{4pt}
\begin{tabular}{@{}lrrrrr@{}}
\toprule
Missingness & Intercept & Complete & $L_1$ only & Reduction & MC interval \\
\midrule
Low & $0.19998$ & 51.16\% & 37.92\% & 1.47\% & 0.68--2.25 \\
Moderate & $-1.09906$ & 24.01\% & 65.06\% & 3.66\% & 2.52--4.85 \\
High & $-2.00020$ & 11.61\% & 77.45\% & 5.05\% & 3.81--6.29 \\
\bottomrule
\end{tabular}
\begin{tablenotes}[flushleft]\footnotesize
\item Reduction is $100\{1-\widehat{\Var}(\widehat\psi_{\seq})/\widehat{\Var}(\widehat\psi_{\cc})\}$, using centered empirical variances. The interval is a paired-bootstrap 95\% Monte Carlo interval from 10,000 resamples; it is expressed in percentage points.
\end{tablenotes}
\end{threeparttable}
\end{table}
\begin{table}[htbp]
\centering
\begin{threeparttable}
\caption{Identification sensitivity with population complete-case proportion fixed at 24\% (5,000 replicates per row).}
\label{tab:sensitivity-identification}
\small
\setlength{\tabcolsep}{4pt}
\begin{tabular}{@{}rrrrrr@{}}
\toprule
$\gamma$ & Intercept & CC bias & CC coverage & Seq. bias & Seq. coverage \\
\midrule
$0.75$ & $-1.47346$ & $-0.01363$ & 83.22\% & $0.00011$ & 95.06\% \\
$1.50$ & $-1.91220$ & $-0.02633$ & 60.88\% & $-0.00000$ & 94.90\% \\
$2.25$ & $-2.41753$ & $-0.03696$ & 44.28\% & $-0.00018$ & 95.56\% \\
\bottomrule
\end{tabular}
\begin{tablenotes}[flushleft]\footnotesize
\item CC denotes the coarsened oracle procedure using $G,Q_0,\rho$; Seq. denotes sequential augmentation. Intercepts solve the stated population response-share equation by deterministic integration.
\end{tablenotes}
\end{threeparttable}
\end{table}

\begin{figure}[htbp]
\centering
\includegraphics[width=0.98\textwidth]{sensitivity_summary.pdf}
\caption{Sensitivity summaries from the same results as \cref{tab:sensitivity-efficiency,tab:sensitivity-identification}. Left: centered empirical variance reduction and paired-bootstrap Monte Carlo intervals. Right: bias and Monte Carlo intervals computed as the estimated bias plus or minus 1.96 Monte Carlo standard errors.}
\label{fig:sensitivity}
\end{figure}

\FloatBarrier
\subsection{A design with bounded treatment probabilities}\label{sec:bounded}

To assess estimated-nuisance behavior under uniform positivity, we also specify a prospective full-data law. Independently, $L_1\sim\mathrm{Bernoulli}(0.5)$ and $L_2\sim\mathrm{Uniform}(-1,1)$. Let
\begin{align*}
A\mid L_1,L_2&\sim\mathrm{Bernoulli}\{\expit(-0.3+0.6L_1+0.7L_2)\},\\
\mu_a(L_1,L_2)&=0.25+0.15a+0.10L_1+0.05L_2,\\
Y\mid A,L_1,L_2&\sim\mathrm{Beta}\{8\mu_A,8(1-\mu_A)\}.
\end{align*}
Potential outcomes may be generated from these conditional beta laws independently of treatment given the confounders, ensuring conditional exchangeability. The ATE is exactly $0.15$, and $\expit(-1)\le e(V)\le\expit(1)$, approximately $[0.269,0.731]$. Outcomes and response probabilities are also bounded as required for the moment and positivity conditions. The two response mechanisms are the same functions of $A,Y,L_1$ as above. This design supplies an instance of the identifying assumptions and uniform positivity; it does not by itself verify the nuisance-rate conditions.

All nuisance functions are estimated with the fixed five-fold procedure in \cref{app:crossfit}. Each configuration uses 1,000 replicates. No tuning parameter is selected from the Monte Carlo performance. The full-data benchmark uses all simulated confounders; each incomplete-data estimator receives only its observed covariates.

\begin{table}[htbp]
\centering
\begin{threeparttable}
\caption{Estimated-nuisance performance with bounded treatment probabilities.}
\label{tab:bounded}
\small
\setlength{\tabcolsep}{4pt}
\begin{tabular}{@{}rlrrrrr@{}}
\toprule
$n$ & Estimator & Bias & SD & Mean SE & RMSE & Coverage \\
\midrule
\multicolumn{7}{l}{\textit{Common validity}} \\
2,500 & Full data & $-0.00003$ & $0.00661$ & $0.00665$ & $0.00661$ & 95.1\% \\
 & Coarsened & $-0.00007$ & $0.00935$ & $0.00958$ & $0.00935$ & 95.9\% \\
 & Sequential & $-0.00001$ & $0.00890$ & $0.00896$ & $0.00890$ & 96.0\% \\
\addlinespace[3pt]
5,000 & Full data & $-0.00013$ & $0.00455$ & $0.00467$ & $0.00455$ & 95.6\% \\
 & Coarsened & $-0.00019$ & $0.00607$ & $0.00618$ & $0.00607$ & 95.4\% \\
 & Sequential & $-0.00027$ & $0.00548$ & $0.00567$ & $0.00548$ & 96.4\% \\
\addlinespace[3pt]
\multicolumn{7}{l}{\textit{Sequential MAR only}} \\
2,500 & Full data & $0.00019$ & $0.00675$ & $0.00665$ & $0.00675$ & 95.1\% \\
 & Coarsened & $0.00437$ & $0.00870$ & $0.00850$ & $0.00973$ & 91.4\% \\
 & Sequential & $0.00000$ & $0.01227$ & $0.01104$ & $0.01227$ & 95.4\% \\
\addlinespace[3pt]
5,000 & Full data & $-0.00001$ & $0.00476$ & $0.00467$ & $0.00475$ & 94.0\% \\
 & Coarsened & $0.00403$ & $0.00574$ & $0.00566$ & $0.00701$ & 87.5\% \\
 & Sequential & $-0.00003$ & $0.00643$ & $0.00632$ & $0.00643$ & 95.6\% \\
\addlinespace[3pt]
\bottomrule
\end{tabular}
\begin{tablenotes}[flushleft]\footnotesize
\item SD is the empirical standard deviation; mean SE is the average estimated standard error. Coverage uses nominal 95\% normal intervals. Each configuration uses 1,000 replicates; the true ATE is 0.15. Coverage Monte Carlo SE is at most 1.58 percentage points (about 0.69 near 95\% coverage).
\end{tablenotes}
\end{threeparttable}
\end{table}
Across the four bounded-design configurations, sequential coverage ranges from 95.4\% to 96.4\%. The bias, empirical variability, and average estimated standard errors are shown together so that coverage can be assessed against its Monte Carlo uncertainty. The common-validity and sequential-only scenarios must still be distinguished when interpreting the coarsened results. A smaller bias does not guarantee a smaller finite-sample mean squared error: in the sequential-only configuration at $n=2{,}500$, sequential RMSE is $0.01227$, compared with $0.00973$ for the biased coarsened estimator. The efficiency-bound comparison does not establish a uniform finite-sample MSE ordering.

\subsection{Estimated nuisances in the stress design}\label{sec:estimated-nuisance}

The corresponding Gaussian-design results appear in \cref{app:crossfit}, including sample sizes 2,500 through 20,000 and oracle substitutions for selected nuisance layers. With all nuisances estimated, sequential coverage ranges from 83.8\% to 89.6\% across the eight sample-size and scenario configurations. The observed undercoverage is a limitation of this estimator and learner in that design; a reduction in bias or improvement with sample size does not establish the asymptotic rate conditions. The oracle-layer comparisons are diagnostics of the specified fitting procedure, rather than a general decomposition of all sources of finite-sample error.

\FloatBarrier
